% Run GCT_KRON_20261120; GPT-6 Astra Pro; effort not exposed.
\section{The exact partial-Hessian boundary}

\subsection*{A-035. Every proper partial bracket survives on the determinant}
Let \(n\geq3\), \(R=n^2\), and \(2\leq r\leq R\). Put
\[
 q_r=r(n-2),\qquad D_r=r+1,\qquad
 \lambda_{n,r}=(n+q_r+1,3,2^{r-2}).
                                                               \tag{A-035.1}
\]
Then
\[
 \boxed{\quad
 K_n(\lambda_{n,r})=a_{\lambda_{n,r}}(D_r[n])=1
                  \qquad(2\leq r<R).
 \quad}                                                        \tag{A-035.2}
\]
At \(r=R\), A-030 instead gives
\(K_n(\lambda_{n,R})=0\) and
\(a_{\lambda_{n,R}}(D_R[n])=1\).
Thus the basic Hessian-bracket obstruction occurs precisely
at the full matrix-variable dimension within this family.
In particular, the proper partial bracket cannot give an
actual permanent multiplicity obstruction: its determinant
multiplicity already equals the complete ambient multiplicity.

For every internally padded permanent with \(1\leq m<n\),
the entire basic bracket family therefore satisfies
\[
       q^m_{\lambda_{n,r}}(D_r[n])
                     \leq K_n(\lambda_{n,r})
                      \qquad(2\leq r\leq n^2).                \tag{A-035.2a}
\]
Indeed, if \(r\leq m^2\), then \(r<n^2\) and the
determinant retains the unique entire ambient copy.
If \(r>m^2\), the length restriction makes the padded
permanent multiplicity zero. This rules out every
member of the basic bracket family at every proper
internal padding, without a large-padding hypothesis.
The separate algebra statement below controls
arbitrary polynomial combinations in its stated
variable range.

For a degree-\(n\) form \(F\), restrict to the first \(r\)
variables and let \(H_r(F)\) be the Hessian determinant
of that restricted form. It has coefficient degree \(r\)
and variable degree \(q_r\). Use the normalized binary
coefficients
\[
 F(xe_1+ye_2)=\sum_i\binom ni F_i x^{n-i}y^i,\qquad
 H_r(F)(xe_1+ye_2)=\sum_i\binom{q_r}i H_i x^{q_r-i}y^i.
\]
The polynomial
\[
                         Z_r=F_0H_1-F_1H_0                 \tag{A-035.3}
\]
has coefficient degree \(D_r\) and highest weight
\(\lambda_{n,r}\). For an upper triangular matrix,
the determinant-line factor is the square of the
determinant of its leading \(r\)-by-\(r\) block.
This is the partial-variable version of the covariance
calculation in A-033.
Taking \(d=r+1\) in F-024 gives
\[
 a_{\lambda_{n,r}}(D_r[n])=1.
                                                               \tag{A-035.4}
\]
Indeed its first part becomes
\((r+1)(n-2)+3=n+q_r+1\), and its length is \(r\),
including the two-row case \(r=2\).
It suffices to prove that \(Z_r\) is nonzero on a
determinant orbit frame.

\paragraph{Up to the smooth-zero Hessian rank.}
Write \(P=\det_n\) and
\[
 A=\operatorname{diag}(1,\ldots,1,0),\qquad B=E_{nn}.
\]
Then \(P(A)=0\) and \(dP_A(B)=1\).
The Hessian bilinear form at \(A\) is
\[
 \operatorname{Hess}P_A(C,D)
 =\operatorname{tr}(C')D_{nn}
       +\operatorname{tr}(D')C_{nn}
       -\sum_{i<n}(C_{in}D_{ni}+D_{in}C_{ni}),
                                                               \tag{A-035.5}
\]
where primes denote the leading \((n-1)\)-by-\((n-1)\)
blocks. This follows by extracting the terms quadratic
in \(u,v\) from \(P(A+uC+vD)\).
Its rank is \(2n\): the trace and \(nn\) coordinates
form one hyperbolic pair, and the \(in,ni\) coordinates
form \(n-1\) further hyperbolic pairs.
In particular the restriction to \(\langle A,B\rangle\)
is nondegenerate.

For any \(2\leq r\leq2n\), enlarge this plane to an
\(r\)-dimensional subspace on which the Hessian is
nondegenerate. Such an extension exists over
\(\mathbb C\), by taking a nondegenerate subspace of
the orthogonal complement in its rank-\(2n\) quotient.
Take its first two basis vectors to be \(A,B\), and
complete this \(r\)-frame to a basis of
\(\operatorname{Mat}_n\).
For the corresponding determinant translate,
\[
             F_0=0,\qquad F_1=1/n,\qquad H_0\ne0.
 \]
Therefore \(Z_r=-H_0/n\ne0\), as required.

\paragraph{Every remaining proper dimension.}
Suppose \(2n<r\leq R\), and put \(k=r-2n\).
The trace pairing on
\(\mathfrak{sl}_{n-1}(\mathbb C)\) is nondegenerate.
Choose a \(k\)-dimensional subspace \(L\) on which
it is nondegenerate, embedded in the leading block
of \(\operatorname{Mat}_n\).
Its maximal permitted dimension is
\[
            (n-1)^2-1=R-2n,
\]
so the choice is possible. Use the \(r\)-dimensional
space with ordered basis beginning
\[
        A,\ B,\ E_{1n},E_{n1},\ldots,E_{n-1,n},E_{n,n-1},
\]
followed by a basis of \(L\). Along
\[
                  X(t)=A+tB=\operatorname{diag}(1,\ldots,1,t)
\]
the Hessian of \(P\) on this space is block diagonal:
\[
 \begin{pmatrix}(n-1)(n-2)t&n-1\\n-1&0\end{pmatrix},
 \quad
 \underbrace{\begin{pmatrix}0&-1\\-1&0\end{pmatrix},
       \ldots,\begin{pmatrix}0&-1\\-1&0\end{pmatrix}}_{n-1},
 \quad -t\,T_L,
                                                               \tag{A-035.6}
\]
where \(T_L=(\operatorname{tr}(C_iC_j))\) is invertible.
For \(t\ne0\), this follows from
\[
 \operatorname{Hess}P_X(C,D)
 =P(X)\{\operatorname{tr}(X^{-1}C)\operatorname{tr}(X^{-1}D)
                    -\operatorname{tr}(X^{-1}CX^{-1}D)\};
\]
polynomiality extends the displayed blocks to \(t=0\).
Their determinant is
\[
                H_r(P)(X(t))=c\,t^k,\qquad
                c=(-1)^{n+k}(n-1)^2\det T_L\ne0.              \tag{A-035.7}
\]
Changing the first basis vector from \(A\) to \(X(t)\)
preserves the determinant of the \(r\)-frame change,
so the same formula evaluates \(H_0\) on the frame
whose first two directions are \(X(t),B\).
Since \(P(X(t))=t\) and \(dP_{X(t)}(B)=1\),
\[
 Z_r=\frac{c k}{q_r}t^k-\frac c n t^k
     =-\frac{2(R-r)}{nq_r}\,c\,t^k.                           \tag{A-035.8}
\]
For every proper \(r<R\), its value at \(t=1\)
is nonzero. This proves (A-035.2).
At \(r=R\), the factor in (A-035.8) is zero,
consistent with the full Hessian identity in A-030;
that identity proves vanishing on every determinant
frame, not merely the frames displayed here.

\paragraph{A whole partial-Hessian algebra cannot separate in the padded range.}
There is a stronger limitation at arbitrary coefficient
degree, beyond the single multiplicity-one statement.
For \(2\leq r\leq2n\), let \(\mathcal J_{n,r}\) be the
space of coefficients of degree-\(n\) forms in
\(x,y,z_3,\ldots,z_r\) whose total \(z\)-degree is at
most two:
\[
 \mathcal J_{n,r}=
 \operatorname{Sym}^n\langle x,y\rangle
 \oplus
 \bigl(\operatorname{Sym}^{n-1}\langle x,y\rangle\otimes Z\bigr)
 \oplus
 \bigl(\operatorname{Sym}^{n-2}\langle x,y\rangle
                                    \otimes\operatorname{Sym}^2Z\bigr).
                                                               \tag{A-035.9}
\]
Here \(Z=\langle z_3,\ldots,z_r\rangle\); when \(r=2\)
the last two summands are absent.
By O-027, determinant matrix pencils project dominantly
onto this entire coefficient space. Equivalently,
restriction to the determinant closure is injective on
the polynomial algebra \(\mathbb C[\mathcal J_{n,r}]\)
embedded in the ambient form-coordinate ring.

Every coefficient of
\[
                 F(x,y,0),\qquad H_r(F)(x,y,0)
 \]
is a polynomial in the coordinates of
\(\mathcal J_{n,r}\). To see this, each entry of
the Hessian matrix along \(z=0\) differentiates at
most two transverse variables. It depends only on
coefficient generators of total transverse degree
zero, one or two. Taking its determinant and
extracting a binary coefficient preserves that
polynomial dependence.
The same statement holds for every Hessian minor
using directions among these first \(r\) variables.
Thus the entire algebra generated by these binary
form coefficients and partial-Hessian coefficients,
including arbitrary polynomial combinations and
all coefficient degrees, restricts injectively to
the determinant closure.

Precisely, if \(V\) is any finite-dimensional space
of ambient highest-weight functions of one degree
and one weight contained in this algebra, then
for any form orbit closure \(\mathcal Y\),
\[
 \dim\operatorname{res}_{\mathcal Y}V
          \leq\dim V
          =\dim\operatorname{res}_{\mathcal D_n}V
          \leq K_n(\operatorname{weight}V).
                                                               \tag{A-035.10}
\]
Consequently its own surviving copies cannot certify
a strict actual permanent multiplicity advantage.
This conclusion does not identify the entire
highest-weight multiplicity space of that type.
Other highest-weight vectors outside the algebra
could change the comparison. If full ambient
coverage is separately proved, then and only then
the injection gives full multiplicity equality;
O-027 supplies that coverage when
\(\ell(\lambda)\leq2n\) and \(\lambda_3\leq2\).

An internally padded permanent uses at most \(m^2\)
variables. A useful partial Hessian must therefore
use at most \(r\leq m^2\) variables. Whenever
\[
                              m^2\leq2n,
                                                               \tag{A-035.11}
\]
every such partial-Hessian binary coefficient algebra
is covered by (A-035.10).
In particular, the complete binary-pair strategy that
proved A-034 cannot transfer directly to the large-padding
range through smaller Hessians.
This closes that proposed construction route; it
neither proves the general padded inequality \(q\leq K\)
nor excludes all higher multiplicities of its shapes.
The stretches of a proper partial-bracket shape may
contain ambient highest-weight functions involving
transverse coefficients of degree at least three,
which are outside the algebra tested here.

\paragraph{Geometric dictionary and scope.}
The restriction, Hessian, binary-coefficient and
evaluation maps are actual polynomial maps over
\(\mathbb Q\). The determinant frame computation
detects the unique entire ambient copy supplied by
F-024; the algebra injection instead uses the full
coefficient dominance proved in O-027.
These are different reasons for passing from
polynomial constructions to multiplicities.
Their comodule and weight interpretations are SGA 3,
I, Propositions 4.7.2--4.7.3, pp.~34--36; their
Schur interpretation is SGA 6, VI, Theorem 3.3,
pp.~392--393, and Macdonald, I, Appendix A(5.3)--(5.6),
pp.~155--156.
All irreducible and Hessian-nondegeneracy assertions
are made in characteristic zero. No positivity of
a virtual class or transfer from an individual
equation to a full multiplicity comparison is used.
